\section{Introduction}

Joint-embedding predictive architectures (JEPAs) offer an attractive route to
visual world models: encode observations, predict future embeddings under
candidate actions, and plan using distances in the learned latent space.
LeWorldModel (LeWM) demonstrates that this recipe can be trained end-to-end
from pixels with a compact prediction objective and a statistical
anti-collapse regularizer \citep{maes2026leworldmodel}.  Yet preventing global
collapse does not ensure that the finite latent budget is allocated to the
factors needed by a planner.

The Obsessed Encoder study makes this distinction concrete
\citep{enigma2026obsessed}.  A tiny, low-entropy, temporally predictable visual
tag can occupy a disproportionate share of a JEPA representation.  The model
continues to satisfy its anti-collapse objective and can achieve low latent
prediction error, but its planner compares observations largely by the tag
rather than by task content.  Filtering static information is not a general
solution: a task goal can itself be static.  The relevant distinction is not
dynamic versus static, but whether changing a factor changes future decisions.

We begin from a simple observation.  The latent prediction loss has two
separable effects.  It must train the predictor to map a current latent and an
action to a future latent.  At the same time, its gradient through the target
and context encoders determines what the representation becomes.  These roles
need not receive identical optimization access.  We therefore preserve the
full predictive learning signal for the predictor, but treat its encoder-side
gradient as a proposal whose priority is informed by an action-derived control
objective.

This intervention is deliberately narrower than claiming that action
supervision identifies every useful physical variable.  Inverse dynamics can
make action-relevant information decodable without ensuring that it dominates
the planner-facing metric.  Conversely, a direction orthogonal to the current
inverse-dynamics gradient may encode already learned control information,
long-horizon consequences, or a goal variable absent from the short-horizon
proxy.  Our central question is therefore:

\begin{quote}
How can a predictive world model reduce nuisance takeover while retaining a
reliable learning path for dynamics not yet covered by its control guide?
\end{quote}

Our contributions are:
\begin{itemize}
    \item We separate \emph{predictor fitting} from \emph{encoder allocation}
    and diagnose predictable-nuisance failure at their optimization interface.
    \item We introduce an asymmetric, per-sample rule that preserves full
    predictor learning while admitting control-supported prediction gradients
    into the encoder and attenuating unsupported components.
    \item We show on watermarked PushT that planning failure can persist even
    when physical variables are linearly decodable, whereas changing their
    relative geometric prominence restores planning.
    \item We introduce falsifying ablations---negative-conflict-only routing,
    norm-matched scalar attenuation, shuffled guides, soft residual admission,
    and component-gradient span diagnostics---to determine whether gains arise
    from direction selection or merely reduced prediction-gradient magnitude.
    \item We evaluate the boundary between nuisance suppression and dynamics
    preservation across five task settings without using privileged physical
    state during training.
\end{itemize}

We do not claim to invent gradient surgery.  Our contribution is the failure
mechanism, the asymmetric JEPA encoder--predictor interface, and an empirical
account of which parts of predictive learning should be allowed to organize a
planner-facing representation.
